% Energy calculation and battery life


% --- Battery life vs. duty cycle table ---
\begin{frame}{Battery Life vs.\ Duty Cycle (3500\,mAh)}
\begin{columns}[T,onlytextwidth]
\begin{column}{0.60\textwidth}
\centering
\renewcommand{\arraystretch}{1.5}
{\small
\setlength{\tabcolsep}{8pt}
\begin{tabular}{r r r r}
\toprule
\textbf{D (\%)} & \textbf{Active/hr} & \textbf{$I_{mean}$} & \textbf{Battery life} \\
\midrule
0.01 & 0.4\,s  & $52\,\mu\text{A}$  & 7.7 years \\
0.1  & 3.6\,s  & $70\,\mu\text{A}$  & 5.7 years \\
0.5  & 18\,s   & $150\,\mu\text{A}$ & 2.7 years \\
\rowcolor{sintefgreen!30}
\textbf{1.0} & \textbf{36\,s} & $\mathbf{250\,\mu\text{A}}$ & \textbf{1.6 years} \\
5.0  & 3\,min  & $1.05\,\text{mA}$  & 4.6 months \\
10.0 & 6\,min  & $2.05\,\text{mA}$  & 2.3 months \\
20.0 & 12\,min & $4.04\,\text{mA}$  & 1.2 months \\
50.0 & 30\,min & $10.0\,\text{mA}$  & 14.5 days \\
\bottomrule
\end{tabular}
}
\vspace{0.3cm}
{\small $I_{mean} = D \cdot 20\,\text{mA} + (1-D) \cdot 50\,\mu\text{A}$}
\end{column}
\begin{column}{0.36\textwidth}
  \begin{block}{$D = 1\%$: operating point}
  \small 36\,s active per 3600\,s; ADXL362 INT1 gates all MCU wake events.
  \end{block}
  \vspace{0.4cm}
  \begin{itemize}\setlength\itemsep{8pt}
  \small
    \item ADXL362 activity threshold: MCU wakes only when acceleration exceeds configured mg limit
    \item ESP-NOW TX: $1.5$\,ms at ${\sim}200$\,mA peak; contribution $<0.01$\,mA to $I_{avg}$
    \item $t_{life} \approx \mathbf{1.6}$\,years on 3500\,mAh
  \end{itemize}
\end{column}
\end{columns}
\note{
  The key to achieving 1% duty cycle is the two-stage hardware gate: the ADXL362
  activity detector only fires INT1 when accelerations exceed the configured threshold,
  so ambient vibrations (traffic, footsteps, HVAC) never wake the MCU.
  A naive design sampling on a timer at 10% duty cycle would drain the battery in
  about 2.3 months on the same battery. The hardware gate buys roughly 8x more
  battery life at our target alarm rate.
}
\end{frame}

% --- Battery life vs. duty cycle table (DEVIATION) ---
\begin{frame}[shrink=15]{Battery Life vs.\ Duty Cycle (3500\,mAh) DEVIATION}
\begin{columns}[T,onlytextwidth]
\begin{column}{0.60\textwidth}
\centering
\renewcommand{\arraystretch}{1.5}
{\small
\setlength{\tabcolsep}{8pt}
\begin{tabular}{r r r r}
\toprule
\textbf{D (\%)} & \textbf{Active/hr} & \textbf{$I_{mean}$} & \textbf{Battery life} \\
\midrule
0.01 & 0.4\,s  & $55\,\mu\text{A}$  & 7.3 years \\
0.1  & 3.6\,s  & $100\,\mu\text{A}$ & 4.0 years \\
0.5  & 18\,s   & $300\,\mu\text{A}$ & 1.3 years \\
\rowcolor{sintefred!30}
\textbf{1.0} & \textbf{36\,s} & $\mathbf{550\,\mu\text{A}}$ & \textbf{8.7 months} \\
5.0  & 3\,min  & $2.55\,\text{mA}$  & 1.9 months \\
10.0 & 6\,min  & $5.05\,\text{mA}$  & 29 days \\
20.0 & 12\,min & $10.0\,\text{mA}$  & 14.5 days \\
50.0 & 30\,min & $25.0\,\text{mA}$  & 5.8 days \\
\bottomrule
\end{tabular}
}
\end{column}
\begin{column}{0.36\textwidth}
  \begin{block}{$D = 1\%$: Worst-case deviation}
  \small Active time per alarm: $3.0\,\text{s}$\\
  (2.5\,s sampling/FFT + 0.5\,s Wi-Fi retry peak)
  \end{block}
  \vspace{0.4cm}
  \begin{itemize}\setlength\itemsep{8pt}
  \small
    \item \textbf{Active current deviation}: rises to $50\,\text{mA}$ average due to prolonged $200\,\text{mA}$ TX retries
    \item \textbf{Significant reduction}: $t_{life}$ drops from $1.6\,\text{years}$ to $\mathbf{8.7\,\text{months}}$
    \item Highlights the importance of low-power ESP-NOW over regular Wi-Fi
  \end{itemize}
\end{column}
\end{columns}

\vspace{0.5cm}
\begin{columns}[c,onlytextwidth]
\begin{column}{0.54\textwidth}
  \centering
  \normalsize
  $I_{active} = \frac{2.5\,\text{s} \cdot 20\,\text{mA} + 0.5\,\text{s} \cdot 200\,\text{mA}}{3.0\,\text{s}} = \mathbf{50\,\text{mA}}$
\end{column}
\hfill
\begin{column}{0.42\textwidth}
  \centering
  \normalsize
  $I_{mean} = D \cdot I_{active} + (1-D) \cdot 50\,\mu\text{A}$
\end{column}
\end{columns}

\note{
  This worst-case deviation slide shows the impact of a sustained 200 mA peak
  during transmit retries and longer acquisition. The active current rises
  from 20 mA to an average of 50 mA during the active interval. At our target
  1% duty cycle, this drops the battery life from 1.6 years to just 8.7 months,
  underscoring the importance of optimizing both the sleep cycles and the link budget.
}
\end{frame}
